In astronomy, the group of objects that make up temporary phenomena such as supernovae (SNe), gamma ray bursts and tidal disruption events are known as transients. They are one of the most important phenomena to observe, as the data can lead to extraordinary discoveries such as the accelerating expansion of the Universe \parencite{riess_observational_1998, perlmutter_measurements_1999}. SNe are particularly important due to a specific type of supernova (SN) called type Ia Supernovae (SNe Ia). 

When SNe Ia are observed, their brightness can be plotted as a function of time. Following standardisation of the light curve, the light curves all reach a well known peak luminosity with some remaining scatter \parencite{martin_o_2024, kelsey_concerning_2022, ginolin_ztf_2025}. This peak luminosity can be used to then find the distance to the location of the detonation, using these distance calculations, the parameters that describe the properties of the Universe (\(\Omega_m, H_0, w, etc\)) can be determined. 

However, the current methods of standardisation show that systematic errors lie in the standardisation of these SNe. If these systematic errors can be accounted for, then truly representative light curves of these SNe can be used to produce more precise cosmological calculations. One of the methods to find new factors to account for in standardisation, is the study of host galaxy properties and their correlation with SNe Ia populations. If the light curves can be standardised based on factors that are intrinsic to the SNe Ia progenitor, then the better standardised the light curves can be which results in better determinations of \(\Omega_m, H_0, w, etc\).